\section*{第 \ref{ch:b2:nvs} 章 --- 赋范向量空间}

\begin{solution}{exo:b2:nvs:1}
单位球：菱形（$\norm\cdot_1$）、圆盘（$\norm\cdot_2$）、
正方形（$\norm\cdot_\infty$），依次嵌套。不等式：
$\norm x_\infty \leq \norm x_2$（单个平方项不超过和）；
$\norm x_2 \leq \norm x_1$（两边平方：$x_1^2 + x_2^2 \leq (\abs{x_1}
+ \abs{x_2})^2$）；$\norm x_1 \leq 2\norm x_\infty$（两项，每项
$\leq \max$）。取等情形：$(1, 0)$ 使前两个不等式取等号；
$(1, 1)$ 使 $\norm x_1 = 2\norm x_\infty$ 取等，并说明
$\norm x_2 = \sqrt2 \norm x_\infty$ 与 $\norm x_1 = \sqrt2 \norm
x_2$ 为另一方向上的极端比值。
\end{solution}

\begin{solution}{exo:b2:nvs:2}
范数公理：齐次性与三角不等式
逐项成立；分离性：$N(f) = 0$ 推出 $f' = 0$（故 $f$ 为常数）
且 $f(0) = 0$，故 $f = 0$，从而是范数。

比较：$\norm f_\infty \leq N(f)$，因为 $\abs{f(x)} \leq
\abs{f(0)} + \abs{\int_0^x f'} \leq \abs{f(0)} +
\norm{f'}_\infty$。逆命题不成立，取 $f_n(x) = \frac1n
\sin(nx)$，则 $\norm{f_n}_\infty \leq \frac1n \to 0$，而
$N(f_n) = 0 + \norm{\cos(nx)}_\infty = 1$。不存在常数 $C$ 使
$N \leq C\norm\cdot_\infty$ 成立。
\end{solution}

\begin{solution}{exo:b2:nvs:3}
$\abs{u(f)} \leq \norm f_\infty \int_0^1 \eu^t\,\dd t = (\eu -
1)\norm f_\infty$，当 $f \equiv 1$ 时取等号：$\vertiii u =
\eu - 1$。

移位：$\norm{S(x)}_\infty = \max(\abs{x_2}, \dots, \abs{x_n}) \leq
\norm x_\infty$，在 $x = e_2$ 处取等号：$\vertiii S = 1$（当 $n
\geq 2$）。
\end{solution}

\begin{solution}{exo:b2:nvs:4}
上界：对 $\norm x_\infty \leq 1$，
\[
\abs{(Ax)_i} = \Bigl|\sum_j a_{ij}x_j\Bigr| \leq \sum_j
\abs{a_{ij}},
\]
故 $\norm{Ax}_\infty \leq \max_i \sum_j \abs{a_{ij}}$。取等号时，设
$i_0$ 使最大值达到，并取 $x_j = \operatorname{sign}(a_{i_0
j})$（模为 $1$ 的元素），则 $(Ax)_{i_0} = \sum_j \abs{a_{i_0
j}}$。由此得公式。对所给矩阵，行和分别为 $3$ 与
$4$，故 $\vertiii A_\infty = 4$。
\end{solution}

\begin{solution}{exo:b2:nvs:5}
\emph{几何级数：}对 $\vertiii B < 1$，级数 $\sum
B^k$ 在巴拿赫空间 $\mathcal{M}_n(K)$ 中绝对收敛
（见\cref{thm:b2:nvs:absoluteconvergence}，$\vertiii{B^k} \leq
\vertiii B^k$），且
\[
(I - B)\sum_{k=0}^{K} B^k = I - B^{K+1} \longrightarrow I :
\]
由乘积的连续性
（见\cref{prop:b2:nvs:bilinear}），$(I - B)\sum_{k\geq0} B^k = I$，故
$I - B$ 可逆，其逆即为该和（并有
$\vertiii{(I-B)^{-1}} \leq \frac{1}{1 - \vertiii B}$）。

\emph{开性：}设 $A$ 可逆，$\vertiii H <
\frac{1}{\vertiii{A^{-1}}}$，则 $A + H = A(I + A^{-1}H)$，其中
$\vertiii{A^{-1}H} \leq \vertiii{A^{-1}}\vertiii H < 1$，故可逆。
于是以 $A$ 为心的球仍含于 $GL_n$ 中。

\emph{求逆的连续性：}令 $B = -A^{-1}H$，
\[
(A + H)^{-1} - A^{-1} = \bigl((I - B)^{-1} - I\bigr)A^{-1}
= \Bigl(\sum_{k \geq 1} B^k\Bigr) A^{-1},
\]
其范数 $\leq \frac{\vertiii B}{1 - \vertiii B}\vertiii{A^{-1}}
\to 0$，当 $H \to 0$ 时。
\end{solution}

\begin{solution}{exo:b2:nvs:6}
连续 $\Rightarrow$ 核为闭集：
$\{0\}$ 的原像为闭集（见\cref{thm:b2:metric:globalcontinuity}）。

反之，设 $\ker\varphi$ 为闭集且 $\varphi \neq 0$。取
$a$ 使 $\varphi(a) = 1$；由于 $a \notin \ker\varphi$ 且
核为闭集，故存在球 $B(a, r)$ 与之不交。现设 $h \in E$
满足 $\varphi(h) \neq 0$，则向量 $a - \frac{h}{\varphi(h)}$
属于 $\ker\varphi$，从而在 $B(a, r)$ 之外：
\[
\Bigl\Vert \frac{h}{\varphi(h)} \Bigr\Vert \geq r
\quad\Longrightarrow\quad
\abs{\varphi(h)} \leq \frac{\norm h}{r},
\]
当 $\varphi(h) = 0$ 时该式显然同样成立：
由\cref{thm:b2:nvs:continuouslinear}中的界(4)可知其连续。
\end{solution}

\begin{solution}{exo:b2:nvs:7}
设 $f_n$ 的值为 $0$，于 $\intcc{0}{\frac12 - \frac1n}$ 上，其后仿射增至
$1$（于 $\frac12$ 处），随后取 $1$，于 $\intcc{\frac12}{1}$ 上。对 $m
\geq n$，$f_m - f_n$ 的支撑位于长度为
$\frac1n$ 的区间上，取值在 $\intcc{-1}{1}$ 内，故 $\norm{f_m - f_n}_1 \leq
\frac1n$，故为柯西列。

设 $f_n \to f$（按 $\norm\cdot_1$），且 $f$ 连续。在
$\intcc{0}{\frac12 - \delta}$ 上（$\delta$ 固定）：$\int \abs{f} =
\int\abs{f - f_n} \leq \norm{f - f_n}_1 \to 0$，当 $n >
\frac1\delta$ 时，故 $\int_0^{1/2 - \delta}\abs f = 0$，由
严格正性得该处 $f = 0$ --- 对每个 $\delta$ 成立，即 $f = 0$ 于
$\intoo{0}{\frac12}$ 上。类似地，$f = 1$ 于 $\intcc{\frac12}{1}$
（各 $f_n$ 在该处均等于 $1$）。由 $\frac12$ 处的连续性得 $0 =
1$，矛盾。极限不存在：该空间不完备。
\end{solution}

\begin{solution}{exo:b2:nvs:8}
良定且连续：$\abs{\varphi(f)} \leq \sum 2^{-n}
\norm f_\infty = \norm f_\infty$，故 $\varphi$ 是线性泛函
且 $\vertiii\varphi \leq 1$（级数对每个
$f$ 均绝对收敛）。

范数为 $1$：固定 $K$；点 $1, \frac12, \dots, \frac1K$
两两互异，故存在连续函数 $f_K$，$\norm{f_K}_\infty
\leq 1$，满足 $f_K(\frac1n) = (-1)^n$（$n \leq K$），分段
仿射插值，在 $0$ 附近取常值）。于是
\[
\varphi(f_K) \geq \sum_{n=1}^{K} 2^{-n} - \sum_{n > K} 2^{-n}
= 1 - 2^{-K+1} \xrightarrow[K \to \infty]{} 1 .
\]

不可取到：若 $\norm f_\infty \leq 1$ 且 $\varphi(f) = 1$，则每一项
必须取到其最大值：$(-1)^n f(\frac1n) = 1$ 对每个
$n$ 成立（否则某一项的严格不足无法得到补偿，
各项均 $\leq 2^{-n}$）。于是 $f(\frac1n) = (-1)^n$；但
$\frac1n \to 0$，而 $f$ 在 $0$ 处连续，迫使
$(-1)^n$ 收敛，产生矛盾。因此上确界不是
最大值——在有限维中这不可能发生，因为其中的闭单位
球是紧的。
\end{solution}

\begin{solution}{exo:b2:nvs:9}
\emph{有界序列：}有界序列落在某个闭
球 $\overline B(0, R) = R\,\overline B(0,1)$，是紧的（
紧的单位闭球在同胚 $x \mapsto Rx$ 下的像）：
于是在其中抽取收敛子列。

\emph{线性泛函：}若 $\dim E < \infty$，则
$E$ 上的每个线性映射均连续（见\cref{thm:b2:nvs:finitedim}）。反之，
设 $\dim E = \infty$（这一点由\cref{thm:b2:nvs:riesz}，
实际上被紧性假设排除——本题的要点在于
两条性质在一般赋范空间中的蕴含关系）：
取规范化的线性无关
序列 $(e_n)$，将其扩充为代数基（在本阶段不加证明地承认），
并定义 $\varphi(e_n) = n$，$\varphi = 0$ 于其余
基向量上，再线性延拓。于是 $\abs{\varphi(e_n)} = n$，
$\norm{e_n} = 1$，故在单位球上无界、不连续。因此
“所有线性泛函连续”迫使空间为有限维。
\end{solution}

\begin{solution}{exo:b2:nvs:10}
对常值函数 $1$ 用柯西–施瓦茨不等式：$\norm f_1 =
\int_0^1 \abs f\cdot 1 \leq \bigl(\int_0^1
f^2\bigr)^{1/2}\bigl(\int_0^1 1\bigr)^{1/2} = \norm f_2$。又
$\norm f_2^2 = \int f^2 \leq \norm f_\infty^2$。取 $f_n(x) =
x^n$：
\[
\norm{f_n}_1 = \frac1{n+1}, \qquad
\norm{f_n}_2 = \frac1{\sqrt{2n+1}}, \qquad
\norm{f_n}_\infty = 1 .
\]
于是 $\norm{f_n}_2/\norm{f_n}_1 = \frac{n+1}{\sqrt{2n+1}} \to
\infty$，且 $\norm{f_n}_\infty/\norm{f_n}_2 = \sqrt{2n+1} \to
\infty$，因此不存在反向不等式，任何两个范数均不等价。
\end{solution}

\begin{solution}{exo:b2:nvs:11}
\emph{距离的下界：}对 $h \in \ker\varphi$，
$\abs{\varphi(x)} = \abs{\varphi(x - h)} \leq
\vertiii\varphi\,\norm{x - h}$；对 $h$ 取下确界：
$d(x, \ker\varphi) \geq \abs{\varphi(x)}/\vertiii\varphi$。

\emph{上界：}不妨设 $\varphi(x) \neq 0$。给定
$\varepsilon > 0$，取单位向量 $u$，使 $\abs{\varphi(u)} \geq
\vertiii\varphi - \varepsilon > 0$，并令 $h = x -
\frac{\varphi(x)}{\varphi(u)}\,u$，则 $\varphi(h) = 0$，且
\[
\norm{x - h} = \frac{\abs{\varphi(x)}}{\abs{\varphi(u)}}
\leq \frac{\abs{\varphi(x)}}{\vertiii\varphi - \varepsilon}.
\]
令 $\varepsilon \to 0$，得 $d(x, \ker\varphi) \leq
\abs{\varphi(x)}/\vertiii\varphi$；等号成立。

\emph{不可取到：}取\cref{exo:b2:nvs:8}中的 $\varphi$
（$\vertiii\varphi = 1$，不可取到），并任取 $x$，使
$\varphi(x) \neq 0$。若某个 $h \in \ker\varphi$ 使
$\norm{x - h} = \abs{\varphi(x)}$，则单位向量 $v = (x -
h)/\norm{x - h}$ 满足 $\abs{\varphi(v)} =
\abs{\varphi(x)}/\norm{x - h} = 1 = \vertiii\varphi$：于是
算子范数将被取到，矛盾。
\end{solution}

\begin{solution}{exo:b2:nvs:12}
$N_1 \leq N_2$ 源于上确界关于定义域的单调性，故
恒等映射 $(E, N_2) \to (E, N_1)$ 为 $1$-利普希茨映射。取 $P_n(x) =
(x/2)^n$：$N_2(P_n) = 1$（在 $x = 2$ 处取到），而 $N_1(P_n) =
2^{-n}$：若 $N_2 \leq CN_1$ 成立，则 $1 \leq C2^{-n}$ 对
所有 $n$ 成立，这不可能。因此两范数不等价，且
其逆恒等映射是不连续的线性双射。

\emph{不完备性：}令 $S_n = \sum_{k=0}^{n}\frac{X^k}{k!}$。
对 $m > n$，$N_1(S_m - S_n) \leq \sum_{k>n}\frac1{k!} \to 0$：
关于 $N_1$ 为柯西列。若 $S_n \to P$（在 $(E, N_1)$ 中），则
逐点收敛，于是 $P(x) = \lim S_n(x) = \eu^x$ 于 $\intcc{0}{1}$ 上；但
次数为 $d$ 的多项式不能在某区间上等于 $\eu^x$（对
$d + 1$ 次求导后左边为零，而 $\eu^x$ 不然）。极限在
$E$ 中不存在：该空间不完备。

在有限维情形下，这三种现象均不可能出现：所有范数
彼此等价，每个赋范空间都完备，且逆
线性双射的逆是有限维空间上的线性映射，
从而连续（见\cref{thm:b2:nvs:finitedim}）。
\end{solution}

\begin{solution}{pb:b2:nvs:1}
\textbf{1.} 值得考虑的候选元构成 $K = \{f \in F
: \norm{x - f} \leq \norm x\}$：非空（$0 \in K$），闭
（闭区间的原像，在连续函数 $f \mapsto
\norm{x - f}$ 之下，与闭集 $F$ 相交，
由\cref{cor:b2:nvs:closedsubspace}），有界（$\norm f \leq
\norm{f - x} + \norm x \leq 2\norm x$）。在
有限维的 $F$ 中，闭且有界即为紧
（由\cref{thm:b2:nvs:finitedim}）；连续函数 $f
\mapsto \norm{x - f}$ 在 $K$ 上取到下确界，该值与
$F$ 上的下确界相等（任一 $f \notin K$ 给出 $\norm{x -
f} > \norm x \geq \inf$）。

\textbf{2.} 在 $(\R^n,
\norm\cdot_2)$ 中的平行四边形恒等式：$\norm{u + v}^2 + \norm{u - v}^2 = 2\norm u^2 +
2\norm v^2$（展开坐标和的平方）。
设 $u \neq v$ 为单位向量：
\[
\Bigl\Vert\frac{u+v}2\Bigr\Vert^2 = 1 - \frac{\norm{u -
v}^2}{4} < 1 .
\]
不是严格凸的：对 $\norm\cdot_\infty$，取 $u = (1, 1, 0,
\dots)$，$v = (1, -1, 0, \dots)$：均为单位向量，其中点
$(1, 0, \dots)$ 的范数为 $1$；对 $\norm\cdot_1$，取 $u = (1,
0, \dots)$，$v = (0, 1, 0, \dots)$：中点 $(\frac12,
\frac12, 0, \dots)$ 的范数为 $1$。

\textbf{3.} 设 $d = d(x, F)$。若 $d = 0$：$x \in \overline F =
F$ 且唯一的最佳逼近元为 $x$。若 $d > 0$ 且 $f_1 \neq f_2$
都是最佳逼近元：$u = \frac{x - f_1}{d}$ 与 $v = \frac{x -
f_2}{d}$ 为两个不同的单位向量，于是
\[
\Bigl\Vert x - \frac{f_1 + f_2}2\Bigr\Vert
= d\,\Bigl\Vert\frac{u + v}2\Bigr\Vert < d ,
\]
而 $\frac{f_1 + f_2}2 \in F$，与
$d$ 的定义矛盾。因此最佳逼近元唯一。

\textbf{4.} $\norm{(0,1) - t(1,0)}_\infty = \max(\abs t, 1)
\geq 1$，等号成立当且仅当 $\abs t \leq 1$：最佳逼近元构成
线段 $\{t(1, 0) : t \in \intcc{-1}{1}\}$，所有点距离均为
$1$ --- 唯一性恰因正方形单位球
具有平直的边而失效（第 2 题）。

\textbf{5.} 设 $M = \max f$，$m = \min f$（最值由紧性
保证可取到）。对任一常数 $c$：$\sup\abs{f - c} \geq
\max(M - c,\, c - m) \geq \frac{M - m}2$，最后一个不等式
成立，因为两者的平均值为 $\frac{M-m}2$；两处等号成立
同时迫使 $M - c = c - m$，即\ $c = c^* = \frac{M + m}2$。
反之 $\sup\abs{f - c^*} = \max(M - c^*, c^* - m) =
\frac{M - m}2$。最佳常数唯一。对 $f(x) = x^2$ 在
$\intcc01$ 上有 $c^* = \frac12$，距离为 $\frac12$。

\textbf{6.} 由 $\cos(n{+}1)\theta + \cos(n{-}1)\theta =
2\cos\theta\cos n\theta$ 可知：$T_0 = 1$，
$T_1 = X$，$T_{n+1} = 2XT_n - T_{n-1}$ 所定义的多项式满足 $T_n(\cos\theta)
= \cos n\theta$（归纳法）。唯一性：两个多项式
在 $\intcc{-1}{1}$ 上（有无穷多点）取值相同者必相等。
归纳法再证 $\deg T_n = n$，且首项系数为
$2^{n-1}$（$n \geq 1$，$T_1$ 的系数为 $1 = 2^0$；递推式
使其加倍）。

\textbf{7.} 每个 $x \in \intcc{-1}{1}$ 都可写成 $\cos\theta$，且
$\abs{\cos n\theta} \leq 1$。取 $\eta_k = \cos\frac{k\pi}n$：
有 $T_n(\eta_k) = \cos k\pi = (-1)^k$，且 $1 = \eta_0 > \eta_1 >
\dots > \eta_n = -1$。根：$\cos n\theta = 0$ 当且仅当 $\theta =
\frac{(2k-1)\pi}{2n}$，即 $n$ 个不同的点
$\cos\frac{(2k-1)\pi}{2n}$，又由 $\frac{(k-1)\pi}n <
\frac{(2k-1)\pi}{2n} < \frac{k\pi}n$，故每个根严格位于
相邻两个极值点之间。

\textbf{8.} $T_2 = 2X^2 - 1$, $T_3 = 4X^3 - 3X$, $T_4 = 8X^4 -
8X^2 + 1$。对 $T_3$：$T_3(1) = 1$，$T_3(\tfrac12) = \tfrac12 -
\tfrac32 = -1$，$T_3(-\tfrac12) = 1$，$T_3(-1) = -1$：具有完美的
正负交错。

\textbf{9.} 设 $u_\pm = x \pm \sqrt{x^2 - 1}$，其中 $x \geq 1$：
它们是 $z^2 - 2xz + 1$ 的根，且 $u_+u_- = 1$。数列
$s_n = \frac{u_+^n + u_-^n}2$ 满足 $s_{n+1} = 2x\,s_n -
s_{n-1}$（由该二次方程导出的牛顿型递推），$s_0 = 1$，
$s_1 = x$：与 $n \mapsto
T_n(x)$ 的递推关系与初值相同，故 $s_n = T_n(x)$ 对一切 $n$ 成立。由于 $0 < u_- \leq 1
\leq u_+$，且 $u_+ > 1$（当 $x > 1$ 时）：$T_n(x) \geq
\frac{u_+^n}2 \to \infty$，且 $T_n(x) \sim \frac12\bigl(x +
\sqrt{x^2-1}\bigr)^n$。（对 $x \leq -1$，利用奇偶性 $T_n(-x)
= (-1)^nT_n(x)$，由递推式易见。）

\textbf{10.} 对每个 $\theta$：$T_m\bigl(T_n(\cos\theta)\bigr)
= T_m(\cos n\theta) = \cos mn\theta =
T_{mn}(\cos\theta)$。多项式 $T_m \circ T_n$ 与
$T_{mn}$ 在 $\intcc{-1}{1}$ 上取值相同，故二者相等。

\textbf{11.} 设 $D = Q_n - P$，其次数 $\leq n - 1$（两者
首项相消）。在极值点处：$(-1)^kD(\eta_k) =
2^{1-n} - (-1)^kP(\eta_k) \geq 2^{1-n} - \abs{P(\eta_k)} > 0$，由假设
可知，故 $D$ 在
递减的 $n + 1$ 个点 $\eta_0 > \dots > \eta_n$ 上取符号交替的非零值。由
介值定理，它至少有 $n$ 个不同的根，
各在一个开区间 $\intoo{\eta_{k}}{\eta_{k-1}}$ 内。
次数 $\leq n - 1$ 的非零多项式不能有 $n$ 个根；
而 $D = 0$ 与严格符号矛盾。矛盾：
$\sup\abs P \geq 2^{1-n}$ 对每个首一多项式 $P$（次数为 $n$）成立。

\textbf{12.} 现有 $(-1)^kD(\eta_k) = 2^{1-n} - (-1)^kP(\eta_k)
\geq 0$，因为 $\abs{P} \leq 2^{1-n}$。在每个
$\intcc{\eta_k}{\eta_{k-1}}$（$k = 1, \dots, n$）上，端点处
$D$ 的值符号弱相反，由介值定理
可得闭区间内的零点 $z_k$。若 $z_k$
可两两互异地选取，则 $D \neq 0$，次数 $\leq n-1$，
有 $n$ 个根，矛盾。相邻两区间
只能共享零点 $z_k = z_{k+1} = \eta_k$，
其中 $0 < k < n$ 为内点。此时 $D(\eta_k) = 0$ 表示 $P(\eta_k) =
(-1)^k2^{1-n}$，是 $P$ 在 $\intcc{-1}{1}$ 上的一个极值
在\emph{内点}处取到：$P'(\eta_k) = 0$；
$\eta_k$ 同样是 $T_n$ 的内点极值：$Q_n'(\eta_k) =
0$。于是 $D'(\eta_k) = 0$，故 $\eta_k$ 为重数
$\geq 2$ 的根，正好补偿共享的区间。总之 $D$
计重数至少有 $n$ 个根，次数 $\leq n -
1$，故 $D = 0$，即 $P = Q_n$。切比雪夫极值定理
得证：唯一的首一最佳逼近元为 $2^{1-n}T_n$，其上确界范数
为 $2^{1-n}$。

\textbf{13.} 次数为 $n$ 的首一多项式恰为
$X^n - R$，其中 $R \in \R_{n-1}[X]$，因此
\[
d_\infty\bigl(X^n, \R_{n-1}[X]\bigr) = \min_{P \text{
首一}}\ \sup_{\intcc{-1}{1}}\abs P = 2^{1-n},
\]
且唯一地在 $R^* = X^n - Q_n$ 处取到。作代换 $x = \frac{1+t}2$：
若 $P$ 是次数为 $n$ 的首一多项式，定义在 $\intcc01$ 上，则 $t \mapsto
2^nP\bigl(\frac{1+t}2\bigr)$ 是 $\intcc{-1}{1}$ 上的首一多项式，且
其上确界为 $2^n\sup_{\intcc01}\abs P$。因此
$\sup_{\intcc01}\abs P \geq 2^{-n}\cdot2^{1-n} = 2^{1-2n}$，等号
恰在 $P^*(x) = 2^{-n}Q_n(2x - 1)$ 时成立：在
$\intcc{0}{1}$ 上距离为 $2^{1-2n}$。

\textbf{14.} $\omega$ 为 $n$ 次首一多项式，故
$\sup_{\intcc{-1}{1}}\abs\omega \geq 2^{1-n}$，等号成立当且仅当
$\omega = Q_n = 2^{1-n}T_n$，即节点恰为 $n$
个 $T_n$ 的根：$x_k = \cos\frac{(2k-1)\pi}{2n}$ ---
即\emph{切比雪夫节点}。最小值为 $2^{1-n}$。等距
节点严格更差；插值误差因子
在节点向端点聚集时取得最小。

\textbf{15.} 手算 $n = 2$ 的情形：$x^2$ 取遍 $\intcc01$，故
$\sup_{\intcc{-1}{1}}\abs{x^2 - c} = \max(\abs c, \abs{1 - c})
\geq \frac12$，在 $c = \frac12$ 时取最小，故最小首一
二次多项式为 $x^2 - \frac12 = \frac12(2x^2 - 1) = Q_2$，其值为
$\frac12 = 2^{1-2}$。对 $n = 10$，在 $\intcc{0}{1}$ 上：$2^{1-20} =
2^{-19} \approx 1.9\cdot10^{-6}$。存在某个 $9$ 次多项式，
与 $x^{10}$ 的偏差不超过百万分之二，在整个
$\intcc01$ 上：此尺度下两者图像无法区分
--- $x^{10}$ 在 $0$ 附近十分平坦，使低次项
承担了全部工作。

\textbf{16.} $\norm{g_k}_\infty = g_k(1) = 1$。对 $j > k$，令
$a = 2^k$，$b = 2^j \geq 2a$，并在 $x_0 = 2^{-1/a}$ 处求值
（于是 $x_0^a = \frac12$）：
\[
g_k(x_0) - g_j(x_0) = \frac12 - \Bigl(\frac12\Bigr)^{b/a}
\geq \frac12 - \frac14 = \frac14 .
\]
故 $\norm{g_k - g_j}_\infty \geq \frac14$ 对所有 $j \neq k$ 成立：任何子列都不是
柯西列，也都不收敛。$\bigl(C(\intcc01), \norm\cdot_\infty\bigr)$ 的
闭单位球不是紧的。

\textbf{17.} $F$ 为闭集（\cref{cor:b2:nvs:closedsubspace}）
且为真子空间：取 $y \notin F$，于是 $\delta = d(y, F) > 0$。由
第 1 题可知该距离在某 $f^* \in F$ 处取到。令 $x
= \frac{y - f^*}{\delta}$，为单位向量（$\norm{y - f^*} =
\delta$）。对每个 $g \in F$：
\[
\norm{x - g} = \frac{\norm{y - (f^* + \delta g)}}{\delta}
\geq \frac{\delta}{\delta} = 1 ,
\]
因 $f^* + \delta g \in F$，故 $d(x, F) \geq 1$；又 $d(x,
F) \leq \norm{x - 0} = 1$：恰为 $1$。

\textbf{18.} 在无限维 $E$ 中归纳构造单位向量：
取 $x_1$ 任意；给定 $x_1, \dots, x_k$ 后，
子空间 $F_k = \operatorname{Vect}(x_1, \dots, x_k)$ 为
有限维的，从而为真子空间，由第 17 题可得单位向量
$x_{k+1}$ 满足 $d(x_{k+1}, F_k) = 1$：特别地，
$\norm{x_{k+1} - x_i} \geq 1$（$i \leq k$）。该序列
两两距离 $\geq 1$：单位球中含有一个序列，
该序列无收敛子列，因而单位球不是紧的 ---
里斯定理，其中常数 $1$ 是最优的。

\textbf{19.} 设 $h_n$ 在
$\bigl[\frac1{n+1}, \frac1n\bigr]$ 的每一半上为仿射函数，从 $0$ 升到
中点处的 $1$，再回到 $0$，其余处为零：连续，
$\norm{h_n}_\infty = 1$。对 $n \neq m$，两者支集至多在
公共端点处相交，且此时两者皆为零；在
$h_n$ 的峰值处 $h_m = 0$，从而 $\norm{h_n - h_m}_\infty = 1$ 恰好成立。对
固定 $x > 0$：$h_n(x) = 0$ 只要 $\frac1n < x$，且
$h_n(0) = 0$ 总成立：$h_n \to 0$ 逐点收敛；但
$\norm{h_n - 0}_\infty = 1$：故不是一致收敛。给出一个显式的
“星座”，相互距离为 $1$，位于单位球内。

\textbf{20.} 半径为 $\frac13$ 的球直径 $\leq
\frac23 < 1$，故其至多含 $h_n$ 中的一个（任两者
距离为 $1$）。有限个这样的球至多只能含
无穷多个 $h_n$ 中的有限个，因此不能覆盖单位球。
全有界性 --- 紧度量空间所具有的这一性质，由
\cref{thm:b2:metric:borellebesgue} 的证明可知，它在此处失效得
最为彻底。

\textbf{21.} 若 $Ax = \lambda x$ 且 $x \neq 0$，则
$\abs\lambda\,\norm x = \norm{Ax} \leq \vertiii A\,\norm x$,
故 $\abs\lambda \leq \vertiii A$。对
$\begin{psmallmatrix}1 & -2\\ 3 & 1\end{psmallmatrix}$:
$\vertiii A_\infty = \max(1 + 2,\ 3 + 1) = 4$
（见\cref{exo:b2:nvs:4}），故每个特征值的模 $\leq
4$；事实上 $\chi_A = X^2 - 2X + 7$ 给出 $\lambda = 1 \pm
\iu\sqrt6$ 的模为 $\sqrt7 \approx 2.65$：该界成立，
但不是最优的。

\textbf{22.} $N_P$ 为范数：$N_P(x) = 0$ 推出 $P^{-1}x = 0$，
故 $x = 0$；齐次性与三角不等式
均从 $\norm\cdot_\infty$ 经线性的 $P^{-1}$ 继承而来。算子
范数：记 $y = P^{-1}x$，$D = \mathrm{diag}(\lambda_i)$，
\[
N_P(Ax) = \norm{P^{-1}AP\,y}_\infty = \norm{Dy}_\infty,
\]
故 $\vertiii A_{N_P}$ 为 $\norm\cdot_\infty$ 算子范数，
即 $D$ 的最大绝对行和
（见\cref{exo:b2:nvs:4}）：$\max_i\abs{\lambda_i}$。

\textbf{23.} 若所有 $\abs{\lambda_i} < 1$，则 $N_P(A^kx) \leq
\rho^kN_P(x)$ 其中 $\rho = \max\abs{\lambda_i} < 1$，故 $A^kx
\to 0$ 对每个 $x$ 成立，并且 $A^k \to 0$ 在
$\mathcal{M}_n$ 的任一范数下成立（有限维中所有范数等价，
\cref{thm:b2:nvs:finitedim}；$A^ke_j$ 对每个
$j$ 而言即逐元收敛）。若某个 $\abs{\lambda} \geq 1$
有特征向量 $x$，则 $\norm{A^kx} = \abs\lambda^k\norm x
\not\to 0$。对 $A = \frac14\begin{psmallmatrix}1 & 2\\ 2 &
1\end{psmallmatrix}$：特征值为 $\frac14(1 \pm 2) =
\frac34, -\frac14$，模均 $< 1$，故 $A^k \to 0$。

\textbf{24.} 二者都是有限维空间
$\R_n[X]$ 上的范数，故对每个 $n$ 等价。取
第 13 题中的最小首一多项式 $P^*_n$（在 $\intcc01$ 上），其
$X^n$ 的系数为 $1$，故 $N_c(P^*_n) \geq 1$，而
$\norm{P^*_n}_{\intcc01} = 2^{1-2n}$。因此
\[
C_n \geq \frac{N_c(P^*_n)}{\norm{P^*_n}_{\intcc01}}
\geq 2^{2n-1} .
\]
等价常数随维数急剧增大：在
并集 $\R[X]$ 上没有统一的常数可用，这恰是
\cref{exo:b2:nvs:12} 中所见的不等价性 --- “所有范数
等价”是关于某一固定维数的定理，
而无限维正是它失效之处。

\textbf{25.} 有限维中闭有界集的紧性
保证了最佳逼近的存在（第
1 题），并为在 $n + 1$ 个极值点处的零点计数提供了动力
（第 11--12 题，借助取到的上确界）。严格凸性支配
最佳逼近的唯一性 --- 圆形的球给出唯一的
极小元，平侧面的球则给出由它们构成的线段（第
2--4 题）。相互距离为
$1$ 的单位向量“星座”（第 17--19 题）破坏了单位球的紧性，
并随之破坏了全有界性（第 20 题）。第 24 题
定量刻画了这一崩塌：联系两个范数的常数
在 $\R_n[X]$ 上以 $4^n$ 的速度增长，故一致的比较在
过渡到 $\R[X]$ 时不复存在。巅峰是切比雪夫的极值
定理（第 11--12 题）：唯一的首一极小元
$2^{1-n}T_n$。最佳逼近的现代归宿是
希尔伯特空间，在那里投影定理以完备性
加平行四边形恒等式取代了紧性 ---
并在本科三年级分册中给出严格的证明。
\end{solution}
